\section*{Capítulo \ref{ch:b1:e-ph-diagrams} --- Diagramas E--pH}

\begin{solution}{exo:b1:e-ph-diagrams:1}
Manganês: \ce{Mn} (0) $<$ \ce{Mn^2+}, \ce{Mn(OH)2} ($+$II) $<$ \ce{MnO2}
($+$IV) $<$ \ce{MnO4-} ($+$VII). Cloro: \ce{Cl-} ($-$I) $<$ \ce{Cl2} (0)
$<$ \ce{HClO}, \ce{ClO-} ($+$I) $<$ \ce{ClO3-} ($+$V).
\end{solution}

\begin{solution}{exo:b1:e-ph-diagrams:2}
$-0.059 \times 3/1 = -0.177$; $+0.059 \times 2/2 = +0.059$;
$-0.059 \times 8/5 = -0.094$; $-0.059 \times 2/2 = -0.059$ V por unidade
de pH. Para o par do cobre, os prótons são liberados do lado da forma
reduzida ($m = -2$): o potencial aumenta com o pH.
\end{solution}

\begin{solution}{exo:b1:e-ph-diagrams:3}
\ce{H+}/\ce{H2} passa por 0~V em pH 0 e nunca chega a 0.50~V para
$\mathrm{pH} \geqslant 0$. \ce{O2}/\ce{H2O} chega a 0.50~V em
$\mathrm{pH} = 0.73/0.059 = 12.4$ e a 0~V só em pH 20.8, fora da escala.
A faixa tem \qty{1.23}{V} de largura em qualquer pH, inclusive em pH 7
(de $-0.41$ a $+0.82$~V).
\end{solution}

\begin{solution}{exo:b1:e-ph-diagrams:4}
(1, 1.0~V): \ce{Fe^3+}. (4, 0): a reta \ce{Fe(OH)3}/\ce{Fe^2+} está em
$1.09 - 0.71 = 0.38$~V e o ponto fica abaixo dela, acima de $-0.47$~V:
\ce{Fe^2+}. (10, $-0.4$~V): entre $-0.07 - 0.59 = -0.66$ e $0.28 -
0.59 = -0.31$~V: \ce{Fe(OH)2}. (10, 0.5~V): \ce{Fe(OH)3}.
\end{solution}

\begin{solution}{exo:b1:e-ph-diagrams:5}
\ce{Fe^2+}/\ce{Fe}: $-0.41 + 0.030 \times (-6) = -0.59$~V.
\ce{Fe^3+}/\ce{Fe(OH)3}: $14.00 - \frac13(38.55 - 6) = 3.15$. Os domínios
dos íons dissolvidos (corrosão) aumentam: \ce{Fe^2+} se estende até
$-0.59$~V e \ce{Fe^3+} até pH 3.15; o metal e os hidróxidos perdem
terreno.
\end{solution}

\begin{solution}{exo:b1:e-ph-diagrams:6}
\ce{Fe(OH)3 + 3H+ + e- -> Fe^2+ + 3H2O}: $E = E^\circ + 0.059\log\frac{[\ce{H+}]^3}{[\ce{Fe^2+}]}
= (E^\circ - 0.059\log c) - 0.177\,\mathrm{pH}$. Em pH 1.82, a reta
encontra \ce{Fe^3+}/\ce{Fe^2+} em 0.77~V, logo a constante é $0.77 + 0.177 \times 1.82
= 1.09$~V.
\end{solution}

\begin{solution}{exo:b1:e-ph-diagrams:7}
$E^\circ(\ce{Cu+}/\ce{Cu}) = 0.52 > E^\circ(\ce{Cu^2+}/\ce{Cu+}) = 0.16$:
pela \cref{prop:b1:e-ph-diagrams:disproportionation-criterion}, \ce{Cu+}
sofre \omterm{def:b1:e-ph-diagrams:disproportionation}{desproporcionamento}. Fronteira única: $E = 0.34 + 0.030\log 0.010 =
0.28$~V.
\end{solution}

\begin{solution}{exo:b1:e-ph-diagrams:8}
Em pH 0, o domínio do cobre (abaixo de 0.28~V) se sobrepõe ao da água e
fica acima da reta do hidrogênio: não há reação com o ácido clorídrico
sem ar. Ácido nítrico: \ce{NO3-}/\ce{NO}, a 0.96~V, fica acima do domínio
do cobre, que é oxidado:
\ce{3Cu + 2NO3- + 8H+ -> 3Cu^2+ + 2NO + 4H2O}.
% ledger: dfg:NO3-_ao, dfg:NO_g (E0 computed in figdata/bachelor-1/standard-potentials.py)
\end{solution}

\begin{solution}{exo:b1:e-ph-diagrams:9}
Em pH 3, a reta \ce{O2}/\ce{H2O} está em $1.23 - 0.18 = 1.05$~V, acima de
todo o domínio de \ce{Fe^2+}: o oxigênio oxida o ferro(II). Acima de pH
1.82, o ferro(III) formado precipita como hidróxido castanho-amarelado:
\ce{4Fe^2+ + O2 + 10H2O -> 4Fe(OH)3 + 8H+}.
\end{solution}

\begin{solution}{exo:b1:e-ph-diagrams:10}
Verticais: $14.00 - \frac12(16.38 - 2) = 6.81$; e, a partir de
\ce{Zn(OH)2 + 2OH- <=> Zn(OH)4^2-} ($K = 10^{-1.93}$) com
$[\ce{Zn(OH)4^2-}] = 10^{-2}$, $[\ce{OH-}]^2 = 10^{-0.07}$, pH 13.97.
\ce{Zn^2+}/\ce{Zn}: $-0.76 + 0.030 \times (-2) = -0.82$~V. \ce{Zn(OH)2}/\ce{Zn}:
inclinação $-0.059$, passando por $(6.81, -0.82)$: $E = -0.42 - 0.059\,\mathrm{pH}$.
\ce{Zn(OH)4^2- + 4H+ + 2e- -> Zn + 4H2O}: inclinação $-0.118$, passando por
$(13.97, -1.24)$: $E = 0.41 - 0.118\,\mathrm{pH}$.
\end{solution}

\begin{solution}{exo:b1:e-ph-diagrams:11}
O domínio do ferro fica abaixo da reta do hidrogênio em qualquer pH. pH
2: o ferro é oxidado a \ce{Fe^2+} com liberação de hidrogênio
(corrosão). pH 7: a reação ainda é possível, mas a reta do hidrogênio
($-0.41$~V) fica perto de \ce{Fe^2+}/\ce{Fe} e a reação é muito lenta. pH
13: o ferro é oxidado a \ce{Fe(OH)2}, um \omterm{def:b1:e-ph-diagrams:corrosion-domains}{domínio de passivação}: forma-se
uma camada e o prego conserva o brilho. Com oxigênio dissolvido, a reta
\ce{O2}/\ce{H2O}, bem acima, oxida o ferro a ferro(III) em qualquer pH: a
ferrugem, \ce{Fe(OH)3} e seus óxidos desidratados.
\end{solution}

\begin{solution}{exo:b1:e-ph-diagrams:12}
O zinco tem o domínio mais baixo: em contato com o aço, ele é oxidado
primeiro, \ce{2Zn + O2 + 2H2O -> 2Zn(OH)2} em pH 8, e os elétrons que
ele libera reduzem o oxigênio na superfície do aço, que permanece
metálica. O magnésio ($E^\circ = -2.36$~V) fica ainda mais abaixo e
também funciona; o cobre (0.34~V) fica acima do ferro: ligado ao cobre,
o ferro seria o metal oxidado, mais depressa do que sozinho.
\end{solution}

\begin{solution}{pb:b1:e-ph-diagrams:1}
\textbf{1.} $-$I, 0, $+$I, $+$I.
\textbf{2.} \ce{Cl-} embaixo, \ce{Cl2}, depois \ce{HClO} e \ce{ClO-} em
cima.
\textbf{3.} \ce{HClO} e \ce{ClO-}, um \omterm{def:b1:predominant-reaction:bronsted}{par ácido--base}.
\textbf{4.} \ce{Cl2 + 2e- -> 2Cl-}.
\textbf{5.} \ce{2HClO + 2H+ + 2e- -> Cl2 + 2H2O}.
\textbf{6.} \ce{ClO- + H2O + 2e- -> Cl- + 2OH-}.
\textbf{7.} \ce{HClO} abaixo de pH 7.55, \ce{ClO-} acima.
\textbf{8.} Nenhum elétron é trocado; a fronteira é fixada apenas por
$K_a$.
\textbf{9.} $1/(1 + 10^{7.4 - 7.55}) = 0.59$: \qty{59}{\percent}.
\textbf{10.} $1/(1 + 10^{0.45}) = 0.26$: três quartos do cloro ativo
seriam o desinfetante mais fraco, \ce{ClO-}.
\textbf{11.} \ce{ClO-}.
\textbf{12.} $E = 1.396 + 0.0296\log\frac{[\ce{Cl2}]}{[\ce{Cl-}]^2} = 1.396
+ 0.0296\log\frac{0.005}{10^{-4}} = 1.396 + 0.050 = 1.446$~V; não há \ce{H+}
na \omterm{def:b1:nernst:couple}{semirreação}.
\textbf{13.} $E = 1.594 + 0.0296\log\frac{[\ce{HClO}]^2[\ce{H+}]^2}{[\ce{Cl2}]}
= 1.594 + 0.0296\log\frac{10^{-4}}{0.005} - 0.0592\,\mathrm{pH} = 1.594 -
0.050 - 0.0592\,\mathrm{pH}$.
\textbf{14.} Igualando, $0.0592\,\mathrm{pH} = 1.594 - 1.396 - 2 \times 0.0503
= 0.0974$, $\mathrm{pH} = 1.646$.
\textbf{15.} \ce{HClO + H+ + 2e- -> Cl- + H2O}: um próton para dois
elétrons, inclinação $-0.0592/2 = -0.0296$.
\textbf{16.} Somando as duas \omterm{def:b1:nernst:couple}{semirreações} (um elétron por átomo de cloro
em cada uma): $2E^\circ = 1.594 + 1.396$, $E^\circ(\ce{HClO}/\ce{Cl-}) =
1.495$~V; fronteira $E = 1.495 - 0.0296\,\mathrm{pH}$.
\textbf{17.} \ce{ClO- + 2H+ + 2e- -> Cl- + H2O}: $E = 1.718 - 0.0592\,\mathrm{pH}$
(constante escolhida por continuidade: $1.495 - 0.0296 \times 7.55 = 1.272 = 1.718
- 0.0592 \times 7.55$).
\textbf{18.} \ce{Cl-} abaixo das retas; \ce{Cl2} em um pequeno triângulo
em $\mathrm{pH} < 1.65$ acima de 1.446~V; \ce{HClO} acima da reta de
inclinação $-0.0296$ até pH 7.55, \ce{ClO-} além disso. A reta
\ce{O2}/\ce{H2O}, de 1.23 a 0.40~V, fica abaixo de todas essas
fronteiras (o diagrama calculado pelo script do capítulo tem a mesma
forma).
\textbf{19.} Não: seu domínio fica acima da reta do oxigênio, de modo que
ele pode oxidar a água, mas a reação é lenta. Uma piscina perde seu cloro
principalmente porque ele oxida o que os banhistas e o ar trazem, e pela
luz do sol; é preciso repô-lo.
\textbf{20.} \ce{Cl2} não tem domínio ali: sofre \omterm{def:b1:e-ph-diagrams:disproportionation}{desproporcionamento},
\ce{Cl2 + H2O -> HClO + Cl- + H+}; em meio básico,
\ce{Cl2 + 2OH- -> ClO- + Cl- + H2O}.
\textbf{21.} Fazendo passar cloro por uma solução de hidróxido de sódio
(a reação da questão 20).
\textbf{22.} \ce{HClO + Cl- + H+ -> Cl2 + H2O}, um \omterm{def:b1:e-ph-diagrams:disproportionation}{comproporcionamento}.
\textbf{23.} Abaixo de pH 1.6, o hipoclorito e o cloreto se
comproporcionam: forma-se cloro, que, além de sua \omterm{def:b1:precipitation:solubility}{solubilidade}, escapa
como um gás tóxico.
\textbf{24.} Em $c = \qty{0.010}{mol/L}$, pH 4 fica fora do domínio de
\ce{Cl2}. Mas a fronteira depende de $c$: repetindo as questões 12--14
para um $c$ qualquer, obtém-se $\mathrm{pH} = 3.34 + \log(2c)$, que chega
a 3.6 para $c = \qty{1}{mol/L}$, a ordem de grandeza da concentração da
água sanitária. Misturar produtos concentrados perto de pH 4 também não é
seguro.
\textbf{25.} \textbf{pH 1.6}: acima dele, em $c = \qty{0.010}{mol/L}$, o
cloro dissolvido sofre \omterm{def:b1:e-ph-diagrams:disproportionation}{desproporcionamento} em ácido hipocloroso e
cloreto.
\end{solution}
